% =============================================================================
% OUTRAS TÉCNICAS E COMO ESCOLHER (16:45 a 16:55)
% =============================================================================
\section[Outras]{Outras técnicas e como escolher}

\begin{frame}{Variáveis instrumentais (IV)}
  \begin{columns}[T]
    \begin{column}{0.5\textwidth}
      \begin{tikzpicture}[scale=0.8,every node/.style={transform shape}]
        \node[cxi] (z) at (0,0) {Instrumento\\$Z$};
        \node[cx] (d) at (3,0) {Participação\\$D$};
        \node[cx] (y) at (6,0) {Resultado\\$Y$};
        \node[cxl] (u) at (4.5,1.8) {Não observados\\$U$};
        \draw[seta] (z) -- (d); \draw[seta] (d) -- (y);
        \draw[setal] (u) -- (d); \draw[setal] (u) -- (y);
        \draw[seta,dashed,red] (z) to[bend right=30] node[below,font=\scriptsize]{sem efeito direto} (y);
      \end{tikzpicture}
    \end{column}
    \begin{column}{0.5\textwidth}
      \footnotesize
      \begin{itemize}
        \item Quando a participação é \alert{endógena} (depende de $U$), usamos algo que a \emph{empurra} sem afetar o resultado diretamente;
        \item Ex.: o \textbf{sorteio} como instrumento da participação (foi o nosso ITT / adesão), distância ao posto, regras de elegibilidade;
        \item Estima o efeito para quem \textbf{responde ao instrumento} (efeito local).
      \end{itemize}
    \end{column}
  \end{columns}
  \vspace{0.3cm}
  \small Pressupostos: instrumento \ind{relevante} (move a participação) e \ind{exógeno} (só afeta $Y$ via $D$).
\end{frame}

\begin{frame}{Controle sintético}
  \begin{itemize}
    \item Para \textbf{uma ou poucas unidades tratadas} (um estado, um país), não há um bom controle único;
    \item Constrói-se um \ind{controle sintético}: média \emph{ponderada} de unidades não tratadas que reproduz a trajetória do tratado \textbf{antes} da política;
    \item O impacto é a \destaque{distância} entre o tratado e o seu sintético \textbf{depois} da política;
    \item Ex. clássico: lei antitabaco da Califórnia (Abadie, Diamond e Hainmueller, 2010).
  \end{itemize}
  \vspace{0.2cm}
  \begin{block}{Relação com o que vimos}
    \small É uma generalização do DiD: em vez de um controle com tendência paralela, \emph{fabrica-se} um controle com peso para cada unidade.
  \end{block}
\end{frame}

\begin{frame}{Como escolher a técnica?}
  \centering
  \begin{tikzpicture}[scale=0.8,every node/.style={transform shape},
      dec/.style={draw=laranja,fill=laranjaclaro,rounded corners=3pt,thick,align=center,inner sep=3pt,font=\scriptsize,minimum width=4.6cm},
      res/.style={draw=azul,fill=azulclaro,rounded corners=3pt,thick,align=center,inner sep=3pt,font=\scriptsize\bfseries,minimum width=3.2cm}]
    \node[dec] (d1) at (0,3.2) {Dá para \textbf{sortear} antes da política?};
    \node[dec] (d2) at (0,1.6) {Há \textbf{corte} de elegibilidade\\num índice contínuo?};
    \node[dec] (d3) at (0,0) {Há um \textbf{instrumento} válido?};
    \node[dec] (d4) at (0,-1.6) {Há dados \textbf{antes e depois}\\de tratados e controles?};
    \node[dec] (d5) at (0,-3.2) {Os fatores da participação\\são \textbf{observáveis}?};
    \node[res] (r1) at (6.2,3.2) {Experimento};
    \node[res] (r2) at (6.2,1.6) {Reg. descontínua};
    \node[res] (r3) at (6.2,0) {Variáveis instrumentais};
    \node[res] (r4) at (6.2,-1.6) {DiD (ou controle sintético\\se há poucas unidades tratadas)};
    \node[res] (r5) at (6.2,-3.2) {Pareamento (PSM)};
    \foreach \i in {1,2,3,4,5} { \draw[setal] (d\i) -- node[above,font=\tiny]{sim} (r\i); }
    \foreach \a/\b in {1/2,2/3,3/4,4/5} { \draw[seta] (d\a) -- node[right,font=\tiny]{não} (d\b); }
  \end{tikzpicture}
\end{frame}

\begin{frame}{As técnicas lado a lado}
  \scriptsize
  \renewcommand{\arraystretch}{1.25}
  \begin{tabularx}{\textwidth}{@{}p{1.6cm}>{\raggedright\arraybackslash}X>{\raggedright\arraybackslash}X@{}}
    \toprule
    \textbf{Técnica} & \textbf{O contrafactual vem de\ldots} & \textbf{Pressuposto-chave} \\ \midrule
    Experimento & Sorteio entre elegíveis & Sorteio bem executado; sem atrito seletivo \\
    Pareamento & Não participantes parecidos & Sem confundidores não observados; região comum \\
    DiD & Tendência do grupo de controle & Tendências paralelas \\
    RDD & Quem está do outro lado do corte & Sem manipulação do corte; efeito local \\
    IV & Variação induzida pelo instrumento & Relevância e exogeneidade do instrumento \\
    Sintético & Combinação ponderada de controles & Bom ajuste pré-tratamento \\ \bottomrule
  \end{tabularx}
\end{frame}
